\documentclass[a4paper,12pt,twoside,openright,titlepage]{book}

%Additional packages
\usepackage[utf8]{inputenc}
\usepackage[T1]{fontenc}
\usepackage[dutch,english]{babel}
\usepackage{imakeidx}
\usepackage{syntonly}
\usepackage[official]{eurosym}
%\usepackage[graphicx]
\usepackage{graphicx}
\graphicspath{ {./images/} }
\usepackage{float}
\usepackage{xurl}
\usepackage{hyperref}
\hypersetup{colorlinks=true, linkcolor=blue, citecolor=blue, filecolor=blue, urlcolor=blue, pdftitle=, pdfauthor=, pdfsubject=, pdfkeywords=}
\usepackage{tabularx}
\usepackage[table]{xcolor} % Table colors
\usepackage{scrextend}
\addtokomafont{labelinglabel}{\sffamily}
\usepackage{listings}
\usepackage{adjustbox}
\usepackage{color}

% Define colors
\definecolor{ashgrey}{rgb}{0.7, 0.75, 0.71}

% Listing style
\lstset{
  backgroundcolor=\color{ashgrey}, % choose the background color; you must add \usepackage{color} or \usepackage{xcolor}; should come as last argument
  basicstyle=\footnotesize,        % the size of the fonts that are used for the code
  breakatwhitespace=true,          % sets if automatic breaks should only happen at whitespace
  breaklines=true,                 % sets automatic line breaking
  extendedchars=true,              % lets you use non-ASCII characters; for 8-bits encodings only, does not work with UTF-8
  frame=single,	                   % adds a frame around the code
  rulecolor=\color{black},         % if not set, the frame-color may be changed on line-breaks within not-black text (e.g. comments (green here))
  keepspaces=true,                 % keeps spaces in text, useful for keeping indentation of code (possibly needs columns=flexible)
  columns=fullflexible,		   % make copy and paste possible
  showstringspaces=false,          % if true show spaces in strings adding particular underscores
  showspaces=false,                % if true show spaces everywhere adding particular underscores; it does not override 'showstringspaces'
}

% Uncomment for production
% \syntaxonly

% Style
\pagestyle{headings}

% Turn on indexing
\makeindex[intoc]

% Define document
\author{D. Leeuw}
\title{Binair}
\date{\today\\v.0.8.0}

\begin{document}
\selectlanguage{dutch}

\maketitle

\copyright\ 2023 Dennis Leeuw\\

\begin{figure}
\includegraphics[width=0.3\textwidth]{CC-BY-SA-NC.png}
\end{figure}

\bigskip

Dit werk is uitgegeven onder de Creative Commons BY-NC-SA Licentie en laat anderen toe het werk te kopi\"eren, distribueren, vertonen, op te voeren, en om afgeleid materiaal te maken, zolang de auteurs en uitgever worden vermeld als maker van het werk, het werk niet commercieel gebruikt wordt en afgeleide werken onder identieke voorwaarden worden verspreid.


%%%%%%%%%%%%%%%%%%%
%%% Introductie %%%
%%%%%%%%%%%%%%%%%%%

\frontmatter
\chapter{Over dit Document}
\section{Leerdoelen}
Dit document geeft een basis uitleg over de boolean algebra in relatie tot de elektrotechniek. De verschillende Boolean functies worden beschreven aan de hand van elektronische schema's. De lezer leert dan ook wat boolean algebra is in de context van de elektrotechniek.


\section{Voorkennis}
\begin{itemize}
\item Om de uitleg te kunnen volgen die beschreven is bij de elektronische werking van de verschillende functies dient de lezer kennis te hebben van de werking van transistoren, weerstanden, spanning en stroom
\item Het is handig als de lezer enigzins vertrouwd is met het lezen van stroomschema's
\end{itemize}


%%%%%%%%%%%%%%%%%
%%% De inhoud %%%
%%%%%%%%%%%%%%%%%
\tableofcontents

\mainmatter
\chapter{Binaire getallen}
Binair is een tweetalig stelsel, vandaar de bi in binair. We tellen dus met alleen de getallen 0 tot 1, dit in tegenstelling tot het decimale stelstel (deca is 10), waarbij we tellen van 0 tot 9.

Laten we bij het vertrouwde decimale stelsel beginnen. We tellen van 0 tot 9 en als er dan nog 1 bij komt (+1) dan zetten we een 1 voor de 0 (10) om aan te geven dat we eenkeer tot 9 hebben geteld. Het is dus tien. Tellen we nu weer door tot 9 (19) en tellen we er dan 1 bij op, dan wordt de voorste 1 een 2 (20) om aan te geven dat we voor de tweede keer tot 9 hebben geteld.

Dit zelfde principe kunnen we ook gebruiken voor binair, alleen tellen we van 0 tot 1 en daarna komt er een 1 voor, dan wordt het dus 10 en daarna 11 waarna we er weer een 1 voor moeten zetten, dus wordt het 100 en zo verder. De getallen 10 en 100 spreken we ook niet uit als tien en honderd, maar als \'e\'en nul en \'e\'en nul nul.


\section{Binair naar decimaal}
Stel we hebben een binair getal 11001110 en we willen weten wat het getal decimaal is. Om binaire getallen om te zetten naar decimaal is er een simpele reken opdracht met machten. Het binaire stelsel kent twee getallen, elke positie in het getal kunnen we weergeven met 2 tot de macht. De macht is dan de waarde van de positie van het bit gerekend vanaf rechts. Om dit wat duidelijker te maken kijken we eerst naar het decimale stelsel, daar ben je wat beter mee vertrouwd.

Het getal honderd is eigenlijk
\begin{math}
	10^2
\end{math}
want
\begin{math}
	10*10=100
\end{math}
zo is 10 eigenlijk
\begin{math}
	10^1
\end{math}
en de getallen onder de 10 zijn eigenlijk
\begin{math}
	10^0
\end{math}
want
\begin{math}
	10^0=1
\end{math}
.

Een getal als 145 kunnen we dus schrijven als
\begin{math}
	1*10^2 + 4*10^1 + 5*10^0
\end{math}
. Deze logica kunnen we ook gebruiken bij het omzetten van binair naar decimaal.


In een binair systeem werken we met een tweetallig stelsel, dus met 2 tot de macht. Ons getal 11001110 kunnen we dus schrijven als:
\begin{math}
	1*2^7 + 1*2^6 + 0*2^5 + 0*2^4 + 1*2^3 + 1*2^2 + 1*2^1 + 0*2^0
\end{math}
Omdat vermenigvuldigen met 0 0 is en vermenigvuldigen met 1 hetzelfde getal oplevert kunnen we het geheel dus nog veel simpeler maken:
\begin{math}
	2^7 + 2^6 + 2^3 + 2^2 + 2^1
\end{math}
. Dit sommetje wordt dus uiteindelijk een simpele optelling:
\begin{math}
	128 + 64 + 8 + 4 + 2 = 206
\end{math}
En daarmee hebben we binair 11001110 naar decimaal 206.


\section{Decimaal naar binair}
Om een decimaal getal om te zetten naar decimaal moeten we van een 10-talig stelsel naar een 2-talig stelsel. Het makkelijkst kunnen we dat doen door te delen door 2. We delen het decimale getal dus steeds door 2 tot we op 0 zijn uitgekomen. Elke keer dat het lukt om door 2 te delen (zonder een breuk te maken) noteren we een 0, als het niet trekken we 1 van het (oneven) getal af en noteren deze 1. Daarna kunnen we het nieuwe evengetal weer delen. De genoteerde nullen en enen levert ons het binaire getal op dat we zoeken. Hieronder volgt een voorbeeld hoe we 1000 decimaal omzetten naar binair:

\begin{tabular} { l r r c }
Bit 1: & 1000 / 2 & 500 & 0 \\
Bit 2: & 500 / 2 & 250 & 0 \\
Bit 3: & 250 / 2 & 125 & 0 \\
Bit 4: & 125 / 2 & 62 & 1 \\
Bit 5: & 62 / 2 & 31 & 0 \\
Bit 6: & 31 / 2 & 30 & 1 \\
Bit 7: & 15 / 2 & 14 & 1 \\
Bit 8: & 7 / 2 & 6 & 1 \\
Bit 9: & 3 / 2 & 2 & 1 \\
Bit 10: & 1 / 2 & 0 & 1 \\
\end{tabular}

Ons binaire getal noteren we van onder naar boven(!) dus als 1111101000.


\section{Bits en bytes}
In de computertechniek spreken we bij binair van een 1 of een 0 als een bit. Zetten we 4 bits als achter elkaar dan hebben we een nibble en 8 bits zijn een byte. Het zijn allemaal woordspelingen op eten. Een bit kan vertaald worden als een stukje of een hapje. To nibble is knabbelen en to bite is bijten/afbijten. Omdat bite te veel leek op bit werd het byte. Het woord nibble voor 4-bits komt niet zo veel meer voor.

Een andere naam voor een byte is een octet van het Griekse octo wat 8 betekent.

Om de woorden bits en bytes te kunnen afkorten is er afgesproken dat de bit een b is en de byte een B, dus de grootste hoeveelheid is de grootste letter.


\section{Duizenden of miljoenen}
Volgens de standaard is een hoeveelheid van 1000 een kilo en van een miljoen een mega. Denk hierbij aan bijvoorbeeld een kilogram. Analoog daaraan is een kb 1000 bits en is kB een 1000 bytes.

Een duizendtal binair is eigenlijk 1024. Voorheen leverde dit verwarringen op. Een hardeschijf van 1024 bits is wat anders dan een hardeschijf van 1000 bits. Om dit onderscheid duidelijk te maken is er voor het binaire 1024 de benaming Kibi bedacht inplaats van kilo. Afgekort wordt het KiB voor 1024 bytes.

\begin{tabular} { l l l l }
Kilo  & k & Kibi & Ki \\
Mega  & M & Mebi & Mi \\
Giga  & G & Gibi & Gi \\
Peta  & P & Pebi & Pi \\
Exa   & E & Exbi & Ei \\
Zetta & Z & Zebi & Zi \\
\end{tabular}


\chapter{Binair rekenen}
\section{Optellen}
Als we twee getallen willen optellen zullen we dus steeds rekening moeten houden dat we heel snel door moeten schuiven met een 1 ervoor. Als we simpel beginnen dan komen we eerst tegen het optellen van twee nullen:
\[ 0 + 0 = 0 \]
Dat is niet anders dan in het decimale stelsel. Ook het optellen van 0 en 1 is niet anders:

\[ 0 + 1 = 1 \]
\[ 1 + 0 = 1 \]

Voor het optellen van twee enen verandert de zaak wel:

\[ 1 + 1 = 10 \]

Nu we dit weten kijk dan eens of de volgende optelling klopt:

\begin{center}
\begin{tabular}{ r l }
	11001110 & \\
	01000101 & +\\
	\hline
	100010011 & \\
\end{tabular}
\end{center}


\section{Aftrekken}
Als we twee getallen van elkaar willen aftrekken dan is ook weer het aftrekken van twee nullen heel simpel:

\[ 0 - 0 = 0 \]

Voor het aftrekken van een nul van \'e\'en is het ook nog eenvoudig:

\[ 1 - 0 = 1 \]

Het wordt moeilijk bij het aftrekken van \'e\'en van nul. We zouden dan uit moeten komen bij -1, maar dat kan in binair niet. Het binaire stelsel kent geen getallen onder nul. We kunnen dus alleen kleine getallen van grote getallen aftrekken zolang de uitkomst niet lager wordt dan 0. We kunnen dus wel 1 aftrekken van 10:

\[ 10 - 01 = 01 \]

Voor het aftrekken van twee enen is de zaak weer simpel:

\[ 1 - 1 = 0 \]

Nu we dit weten kijk dan eens of de volgende aftrekking klopt:

\begin{center}
\begin{tabular}{ r l }
	11001110 & \\
	01000101 & -\\
	\hline
        10001001 & \\
\end{tabular}
\end{center}


\section{Vermenigvuldigen}
Het vermenigvuldigen van binaire getallen is niets relatief simpel, vandaar dat computers ook zulke goede rekenaars zijn. Vermenigvuldigen is schuiven en optellen. Het is waarschijnlijk het makkelijkst met een voorbeeld:

\begin{center}
\begin{tabular}{ r l }
      11001110 & \\
      01000101 & x\\
\hline
      11001110 & \\
     000000000 & \\
    1100111000 & \\
   00000000000 & \\
  000000000000 & \\
 0000000000000 & \\
11001110000000 & + \\
\hline
11011110000110
\end{tabular}
\end{center}

Voor elk bit uit het getal waarmee we vermenigvuldigen schuift het getal dat vermenigvuldigt wordt op naar links en krijgt er een 0 achter. Heeft het getal waarmee we vermenigvuldigen een 1 staan dan kom het opgeschoven getal er te staan, staat er een 0 van is de opteller allemaal nullen, kortom er gebeurt niets.


%%%%%%%%%%%%%%%%%%%%%
%%% Index and End %%%
%%%%%%%%%%%%%%%%%%%%%
\backmatter
\printindex
\end{document}

%%% Last line %%%
